\documentclass{article}
\usepackage{import}
\import{../../template/}{article-template}
\graphicspath{{../../images}}
\addbibresource{../../template/library.bib}

\title{Sloshing}
\author{PENG}
\date{Created: 20260722, Version: 20260916}

\begin{document}
\maketitle
\tableofcontents

\section{Introduction}
\begin{definition}[dispersion]~
    \begin{itemize}
        \item frequency dispersion
        \item amplitude dispersion
    \end{itemize}
\end{definition}

\begin{definition}[nonlinear resonance]
    Nonlinear resonance is a mechanism by which energy is continuously exchanged between a small number of linear wave modes \cite{DUREY_2023_RESONANT_TRIAD_INTERACTIONS_OF_GRAVITY_WAVES_IN_CYLINDRICAL_BASINS}.
\end{definition}

\section{Analytical Method}
\begin{assumption}
    Potential flow:
    \begin{itemize}
        \item irrotational
        \item inviscid
        \item incompressible
    \end{itemize}
\end{assumption}

\begin{definition}
    Governing equations:
    \begin{alignat}{3}
        \label{eq:continuity-equation}
        \Delta\phi + \phi_{zz}                                                             & = 0 \quad        &  & \text{for } x \in D, \quad          &  & -h<z<\epsilon\eta \\
        \label{eq:dynamic-equation}
        \phi_{t} + \eta + \frac{\epsilon}{2}\left( |\nabla\phi|^{2} + \phi_{z}^{2} \right) & = 0 \quad        &  & \text{for } x \in D, \quad          &  & z = \epsilon\eta  \\
        \label{eq:kinematic-equation}
        \eta_{t} + \epsilon\nabla\phi\cdot\nabla\eta                                       & = \phi_{z} \quad &  & \text{for } x \in D, \quad          &  & z = \epsilon\eta  \\
        \label{eq:no-flux-vertial-walls}
        n\cdot\nabla\phi                                                                   & = 0 \quad        &  & \text{for } x \in \partial D, \quad &  & -h<z<\epsilon\eta \\
        \label{eq:no-flux-horizontal-base}
        \phi_{z}                                                                           & = 0 \quad        &  & \text{for } x \in D, \quad          &  & z=-h
    \end{alignat}
\end{definition}

\emph{Why does heave excitation lead to either symmetrical or asymmetrical modes?}

\begin{figure}[h]
    \centering
    \includegraphics{faraday-wave-modes.png}
    \caption{Faraday wave modes}
    \label{fig:faraday-wave-modes}
\end{figure}

\printbibliography
\end{document}
